\subsection{Experiment design and numerical reliability}
The experiments follow the IVP framework of Lecture 10, the truncation
error convention of Lecture 11, and the stability and stiffness analysis
of Lectures 12 and 13~\cite{ae370lecture10,ae370lecture11,ae370lecture12,ae370lecture13}.
Lecture 3 motivates stating the norm used for every error comparison~\cite{ae370lecture3}. These lecture connections describe the methodology.
They do not establish that the pending spacecraft simulations are accurate.

Write the seven component IVP as $\dot{\mathbf y}=f(\mathbf y,t)$ with
initial data at one time. The implementation interface accepts time first.
Classical RK4 is the planned explicit four stage method. With
$h=\Delta t$, its stages evaluate $f$ at the initial time, twice at the
midpoint, and at the final time. The update weights are
$(1,2,2,1)/6$. Magnetic forcing and feedback must be reevaluated at each
stage. The integrator implementation and full derivation belong to the
numerical methods contribution.

Lecture 11 divides the exact starting value step defect by $h$ when defining
$\tau_k$. Under the smoothness assumptions for classical RK4, the unscaled
step defect is $O(h^5)$, this convention gives $\tau_k=O(h^4)$, and stable
integration over a fixed interval gives global error $O(h^4)$. The smooth
verification case must be tested separately from trajectories crossing the
actuator saturation boundary. Such crossings can invalidate a smooth
Taylor expansion across a step.

For example, on $\dot u=\lambda u$, direct substitution of the four stages
gives
\[
 R(z)=1+z+\frac{z^2}{2}+\frac{z^3}{6}+\frac{z^4}{24},
 \qquad z=h\lambda.
\]
The expansion of $e^z-R(z)$ starts at $z^5/120$, which verifies the step
defect order on this scalar test. This calculation alone is not a general
nonlinear order proof. For a frozen decaying linear mode, absolute stability
requires $|R(h\lambda)|<1$. The negative real boundary is approximately
$h|\lambda|=2.785$. The Jacobian spectrum can change along a nonlinear
trajectory. Neutral quaternion modes and growing modes must be recorded
separately. A ratio of nonzero decaying eigenvalue magnitudes is a local
indicator of separated time scales, not proof of stiffness or nonlinear
stability. Production step selection also requires trajectory and metric
refinement, especially near saturation switches.

For paired histories sampled at common nodes, let $d_{q,k}$ be the smaller
of $\|\widehat q_k-\widehat q^{\rm ref}_k\|_2$ and
$\|\widehat q_k+\widehat q^{\rm ref}_k\|_2$, where the hats denote locally
normalized copies. This removes the equivalent quaternion sign choice.
Set $\omega_*=\omega_{\rm threshold}$ and define
\[
 D_k^2=d_{q,k}^2+
 \frac{\|\boldsymbol\omega_k-\boldsymbol\omega_k^{\rm ref}\|_2^2}{\omega_*^2},
 \qquad
 E=\left[\frac{1}{T}\sum_k
 \frac{t_{k+1}-t_k}{2}(D_k^2+D_{k+1}^2)\right]^{1/2}.
\]
The recorded sampled maximum accompanies this dimensionless time RMS.
Raw quaternion norm defects are retained separately. Reference values are
compared only at nested output nodes, without introducing interpolation
error into the state order estimate. For equal refinement ratios,
\[
 p_{\rm obs}=\log_2\left(\frac{E_{h,h/2}}{E_{h/2,h/4}}\right).
\]
An estimate is absent when a difference is zero. Self convergence verifies
neither the physical equations nor the frame convention, so independent
checks remain necessary.

Recovery time is the start of the earliest complete qualifying hold interval
ending by the deadline. Rate threshold equality qualifies. The crossing is
estimated from a linear interpolant of sampled rate norms, an application
of local interpolation from Lecture 4~\cite{ae370lecture4}. The interpolation error for a smooth
scalar history has an $O(h^2)$ bound. At a transverse threshold crossing,
dividing that error by a nonzero crossing slope gives an $O(h^2)$ time error.
Tangencies and unresolved excursions require separate attention. Thus RK4
state accuracy does not imply fourth order recovery time measurements.
Output spacing is refined on the same finest trajectory separately from
integration step refinement. This cannot certify an unresolved excursion
between all retained samples.

The initial pilot uses three spin magnitudes, three gains, three dipole
limits and four common geometries, giving 108 orbital cases. The tentative
production grid has 1152 orbital cases. Attitudes are generated from
normalized Gaussian four vectors, spin directions from normalized Gaussian
three vectors, and orbital phases uniformly on $[0,2\pi)$. The seed is
recorded and geometries are paired across settings. Physical inputs and
derived deadlines come from the validated canonical baseline.

Report recovery times conditional on success, with unsuccessful valid cases
retained as censored. Report simulation errors separately. A missed deadline
is not proof that eventual recovery is impossible. The paired field study
freezes the initial inertial field and its inertial derivative while retaining
the transformation into body components. Parameter choices and the numerical
acceptance tolerances must be agreed before production results are interpreted.
